Read in Gauge Theory Gravity: the metric $g$ below is the derived metric of the gravitational field $(\hb, \Omega)$ on flat spacetime, the manifold is the invertible region, and the tractor connection is the flat-space $\Cl(2,4)$ gauge field of Theorem \claimref{synthesis.gtg-tractor-gauge-field-reproduces-tractor-connection}; the index computation below is its reading in the splitting of $g$ (Fact \claimref{synthesis.imported-gtg-dictionary}).

(i) $de^i = p_i\tau^{p_i-1}d\tau\wedge dx^i = (p_i/\tau)\,e^0\wedge e^i$ and $de^0 = 0$. With $\omega^{i0} = (p_i/\tau)e^i$ and $\omega^{ij} = 0$: $de^i + \omega^{i}{}_{0}\wedge e^0 = (p_i/\tau)(e^0\wedge e^i + e^i\wedge e^0) = 0$ and $de^0 + \omega^{0}{}_{i}\wedge e^i = \sum_i (p_i/\tau)\, e^i\wedge e^i = 0$, so this is the Levi-Civita connection. The tractor connection \eqref{eq:conn} adds the soldering summand, $\mu_b \mapsto g_{ab}\rho$, and the Schouten summand, $\mu_b \mapsto -P_{ab}\sigma$, both constant in the orthonormal frame on the Einstein region, and acts on $I$ and $X$ through them: $\nabla_a X = e_a$ and $\nabla I = 0$, regular. Curvature: $R^{i0} = d\omega^{i0} + \omega^{i}{}_{j}\wedge\omega^{j0} = d\bigl(p_i\tau^{p_i-1}dx^i\bigr) = p_i(p_i-1)\tau^{p_i-2}\,d\tau\wedge dx^i$, and $R^{ij} = \omega^{i}{}_{0}\wedge\omega^{0j} = (p_i/\tau)e^i\wedge(-(p_j/\tau)e^j)$. A gauge pole $D\,d\tau/\tau$ has $d(d\tau/\tau) = 0$ and contributes to the curvature only through $(d\tau/\tau)\wedge[D, \omega_{\mathrm{reg}}]$, a simple pole at most; a double pole with tangential legs is not of that form. (ii) With $R^{ab}{}_{cd}$ read off from (i), the contraction $\mathrm{Ric}_{bd} = -R^{a}{}_{bad}$, which is the one that gives de Sitter $\mathrm{Ric} = +\Lambda g$ in the mostly-minus signature, yields $R_{00} = \sum_i p_i(p_i-1)/\tau^2$ and $R_{ii} = -\bigl(p_i(p_i-1) + p_i\sum_{j\neq i}p_j\bigr)/\tau^2 = -p_i(\sum_j p_j - 1)/\tau^2$, off-diagonal terms vanishing because every $R^{ab}$ is a multiple of $e^a\wedge e^b$. (The opposite contraction negates both displays and would give de Sitter $\mathrm{Ric} = -\Lambda g$, inconsistent with $P_{ab} = (\Lambda/6)g_{ab}$ and with the system of Proposition \claimref{dark-energy.cosmological-constant-does-not-change-residue}, whose de Sitter limit is $3H^2 = \Lambda$.) Vanishing of $R_{ii}$ for some $p_i \neq 0$ gives $\sum p = 1$, and then $R_{00} = 0$ gives $\sum p^2 = 1$. The unit-boost reading is \cite[Prop.~8.26(3)]{ButlerT4} (\claimref{sheaf.matching-fibre}): $B_w$ squares to $\sum p_i^2$ on its boost plane. With $\Lambda$ the corrections are those of Proposition \claimref{dark-energy.cosmological-constant-does-not-change-residue}, of relative order $\Lambda\tau^2$. (iii) Write $R^{0i} = A_i\,e^0\wedge e^i$, $R^{jk} = C_{jk}\,e^j\wedge e^k$ with $A_i = -p_i(p_i-1)/\tau^2$, $C_{jk} = -p_jp_k/\tau^2$. In $\epsilon_{abcd}R^{ab}\wedge R^{cd}$ the pairs $\{a,b\} = \{0,i\}$, $\{c,d\} = \{j,k\}$ and their exchange each contribute four equal terms, so the sum is $8\sum_i A_iC_{jk}\,e^0e^1e^2e^3 = 8\,p_1p_2p_3\sum_i(p_i - 1)\,\tau^{-4}\,\mathrm{vol}$, and $\sum_i(p_i - 1) = 1 - 3 = -2$ on the Kasner relations; with $\mathrm{vol} = \tau^{\sum p_i}d\tau\wedge d^3x = \tau\,d\tau\wedge d^3x$ this is $-16\,p_1p_2p_3\,\tau^{-3}d\tau\wedge d^3x$. (On the leading Kasner term this equals $\sqrt{|g|}$ times the Kretschmann scalar, $64/(27\tau^4)$ for the Schwarzschild exponents (\claimref{dark-energy.imported-weighted-limit-schwarzschild-singularity-matching}), as it must on an Einstein metric.) The Einstein-Cartan density is $-2\sqrt{|g|}\,R$ up to the fixed normalisation, which vanishes on a vacuum and equals $-8\Lambda\tau$ on the Einstein region; the cosmological density is $4!\,\mathrm{vol}$. Exactness of the first term is Theorem \claimref{dark-energy.frozen-point-makes-stelle-west-vector}(2).
